\documentclass[12pt]{article}

\usepackage[a4paper,margin=1in]{geometry}
\usepackage{amsmath,amssymb}
\usepackage{enumitem}
\usepackage{setspace}
\usepackage{titlesec}

\setstretch{1.15}

\titleformat{\section}{\bfseries\large}{\thesection.}{0.5em}{}
\titleformat{\subsection}{\bfseries}{\thesubsection}{0.5em}{}

\begin{document}

\begin{center}
\textbf{CSE400 – Fundamentals of Probability in Computing} \\[4pt]
\textbf{Lecture 5: Bayes’ Theorem, Random Variables, and Probability Mass Function}
\end{center}

\section{Overview}
This lecture covers:
\begin{itemize}[noitemsep]
\item Bayes’ Theorem as a weighted average of conditional probabilities
\item Law of Total Probability and Bayes Formula
\item Concept and motivation of Random Variables
\item Probability Mass Function (PMF) with examples
\end{itemize}

\section{Definitions and Notation}
\begin{itemize}[noitemsep]
\item Events: $A, B, B_1, B_2, \dots, B_n$
\item Complement: $B^c$
\item Probability: $\Pr(\cdot)$
\item Conditional Probability: $\Pr(A \mid B)$
\item Random Variable: $X, Y$
\item Range of $X$: $R_X = \{x_1, x_2, x_3, \dots\}$
\item PMF: $p(x) = \Pr(X = x)$
\end{itemize}

\section{Assumptions and Conditions}
\begin{itemize}[noitemsep]
\item Events $B_1, B_2, \dots, B_n$ are mutually exclusive
\item $\bigcup_{i=1}^n B_i = B$ (or sample space partition)
\item Random variables discussed are primarily discrete
\item Probabilities of all possible outcomes sum to $1$
\end{itemize}

\section{Main Results / Theorems}

\subsection{Weighted Average of Conditional Probabilities}
Let $A$ and $B$ be events. Then:
\[
A = AB \cup AB^c
\]
Since $AB$ and $AB^c$ are mutually exclusive,
\[
\Pr(A) = \Pr(AB) + \Pr(AB^c)
\]
\[
= \Pr(A \mid B)\Pr(B) + \Pr(A \mid B^c)[1 - \Pr(B)]
\]
Hence, $\Pr(A)$ is a weighted average of conditional probabilities.

\subsection{Law of Total Probability (Formula 3.4)}
If $B_1, B_2, \dots, B_n$ are mutually exclusive and exhaustive, then:
\[
\Pr(A) = \sum_{i=1}^n \Pr(A \cap B_i)
= \sum_{i=1}^n \Pr(A \mid B_i)\Pr(B_i)
\]

\subsection{Bayes Formula (Proposition 3.1)}
Using:
\[
\Pr(A \cap B_i) = \Pr(B_i \mid A)\Pr(A)
\]
We obtain:
\[
\Pr(B_i \mid A) =
\frac{\Pr(A \mid B_i)\Pr(B_i)}
{\sum_{j=1}^n \Pr(A \mid B_j)\Pr(B_j)}
\]
Where:
\begin{itemize}[noitemsep]
\item $\Pr(B_i)$ is the \textit{a priori} probability
\item $\Pr(B_i \mid A)$ is the \textit{posteriori} probability
\end{itemize}

\subsection{Random Variables}
A random variable $X$ on sample space $\Omega$ is a function:
\[
X : \Omega \rightarrow \mathbb{R}
\]
It assigns a real number to each sample point.

Key ideas:
\begin{itemize}[noitemsep]
\item Values depend on experiment outcomes
\item Probabilities are assigned to these values
\end{itemize}

\subsection{Distribution of a Random Variable}
For any $a$ in the range of $X$:
\[
\{\omega \in \Omega : X(\omega) = a\}
\]
is an event. \\
$\Pr[X = a]$ defines the distribution of $X$.

\subsection{Probability Mass Function (PMF)}
A random variable is discrete if it takes at most a countable number of values. \\
For discrete $X$ with range $R_X = \{x_1, x_2, \dots\}$, the PMF is:
\[
p(x) = \Pr(X = x)
\]
and:
\[
\sum_k p(x_k) = 1
\]

\section{Proofs / Derivations}

\textbf{Derivation of Bayes Formula}

From Law of Total Probability:
\[
\Pr(A) = \sum_{i=1}^n \Pr(A \mid B_i)\Pr(B_i)
\]
Using conditional probability:
\[
\Pr(B_i \mid A)
= \frac{\Pr(A \cap B_i)}{\Pr(A)}
= \frac{\Pr(A \mid B_i)\Pr(B_i)}
{\sum_{j=1}^n \Pr(A \mid B_j)\Pr(B_j)}
\]

\section{Worked Examples}

\textbf{Example 3.1 (Part 1/2)}

Let:
\begin{itemize}[noitemsep]
\item $A_1$: Policyholder has an accident in 1 year
\item $A$: Policyholder is accident prone
\end{itemize}

Given:
\[
\Pr(A_1 \mid A) = 0.4, \quad
\Pr(A_1 \mid A^c) = 0.2
\]
\[
\Pr(A) = 0.3, \quad \Pr(A^c) = 0.7
\]

\[
\Pr(A_1) = (0.4)(0.3) + (0.2)(0.7) = 0.26
\]

\textbf{Example 3.1 (Part 2/2)}

Find $\Pr(A \mid A_1)$:
\[
\Pr(A \mid A_1) =
\frac{\Pr(A_1 \mid A)\Pr(A)}{\Pr(A_1)}
= \frac{0.4 \times 0.3}{0.26}
= \frac{6}{13}
\]

\textbf{Example 3.2 (Card Problem)}

Cards:
\[
RR, BB, RB
\]
Let $R$ be the event that the upper side is red.
\[
\Pr(RB \mid R) =
\frac{\frac{1}{2} \times \frac{1}{3}}
{(1 \times \frac{1}{3}) + (\frac{1}{2} \times \frac{1}{3}) + (0 \times \frac{1}{3})}
= \frac{1}{3}
\]

\textbf{Random Variable Example}

Toss 3 fair coins. Let $Y$ be the number of heads.
\[
P(Y=0)=\frac{1}{8}, \;
P(Y=1)=\frac{3}{8}, \;
P(Y=2)=\frac{3}{8}, \;
P(Y=3)=\frac{1}{8}
\]
\[
\sum_{i=0}^3 P(Y=i) = 1
\]

\textbf{PMF Example}

Given:
\[
p(i) = c \frac{\lambda^i}{i!}, \quad i = 0,1,2,\dots
\]
Using:
\[
\sum_{i=0}^\infty p(i) = 1
\Rightarrow
c \sum_{i=0}^\infty \frac{\lambda^i}{i!} = 1
\Rightarrow
c e^\lambda = 1
\Rightarrow
c = e^{-\lambda}
\]

\[
P(X=0) = e^{-\lambda}
\]
\[
P(X>2) = 1 - [P(X=0) + P(X=1) + P(X=2)]
\]

\section{Logical Dependencies / Summary}
\begin{itemize}[noitemsep]
\item Law of Total Probability $\rightarrow$ Bayes Formula
\item Bayes Theorem links prior and posterior probabilities
\item Random Variables map outcomes to numbers
\item PMF fully characterizes discrete random variables
\item All probabilities must sum to $1$
\end{itemize}

\end{document}
